% @see : https://coopmaths.fr/alea/?uuid=71e5c&id=2L21-1&n=2&d=10&s=1&s2=1-2-3-4-5-6-7-8-9&cd=1&tip=0&alea=SmM1&uuid=3ad01&id=2L21-5&n=4&d=10&s=0&s2=false&cd=1&tip=0&alea=ML0o&uuid=550cf&id=2L21-6&alea=P6aQ&uuid=bb6d5&id=2L22-2&n=3&d=10&s=1-2-3-5&cd=1&tip=0&alea=Br0N&uuid=31017&id=2L22-4&n=1&d=10&s=3&cd=1&tip=0&alea=XZGz&uuid=27741&id=2L20-2&n=1&d=10&s=1&cd=1&tip=0&alea=ILxS&uuid=622b9&id=2L20-1&n=1&d=10&s=0&s2=false&cd=0&tip=0&alea=lZvf&uuid=cf5b7&id=2L20-4&alea=hcGK&v=eleve
\documentclass[a4paper,11pt,fleqn]{article}

%%% COULEURS %%%
\usepackage[table,svgnames]{xcolor}
\definecolor{nombres}{cmyk}{0,.8,.95,0}
\definecolor{gestion}{cmyk}{.75,1,.11,.12}
\definecolor{gestionbis}{cmyk}{.75,1,.11,.12}
\definecolor{grandeurs}{cmyk}{.02,.44,1,0}
\definecolor{geo}{cmyk}{.62,.1,0,0}
\definecolor{algo}{cmyk}{.69,.02,.36,0}
\definecolor{correction}{cmyk}{.63,.23,.93,.06}
\arrayrulecolor{couleur_theme} % Couleur des filets des tableaux

\usepackage[most]{tcolorbox} % Supprimé à cause de conflits avec ProfCollege
% Découpe des cadres d'exercice. Le style est redéfini par la sortie LaTeX
% juste avant un exercice à garder d'un seul tenant, puis rétabli :
%   \tcbset{mathaleaexo/.style={unbreakable}}
% Passer par un style plutôt que par la clé « breakable » directement évite
% de changer la découpe des autres cadres du document (ProfCollege...).
\tcbset{mathaleaexo/.style={breakable}}
\tcbuselibrary{documentation}

%%%%%%%% Modification paramétrage pour footnotes
% Le paquet semble déjà appelé en amont, curieux qu'il n'y ait pas eu de clash
% avec l'appel ci-dessous mais seulement quand j'ai essayé de l'inclure avec l'option lualatex

% \usepackage{hyperref}
% Parametrages
\hypersetup{
    colorlinks=true,% On active la couleur pour les liens. Couleur par défaut rouge
    linkcolor=black,% On définit la couleur pour les liens internes
    % filecolor=magenta,% On définit la couleur pour les liens vers les fichiers locaux      
    urlcolor=black,% On définit la couleur pour les liens vers des sites web
    % pdftitle={Puissance Quatre},% On définit un titre pour le document pdf
    % pdfpagemode=FullScreen,% On fixe l'affichage par défaut à plein écran
    }
% Pour avoir les numérotations arabic y compris dans les environnements tcolorbox
\renewcommand{\thempfootnote}{\arabic{mpfootnote}}
%%%%%%%%

\usepackage{multicol} 
\usepackage{calc}
\usepackage{enumitem}
\usepackage{graphicx}
\usepackage{tabularx}
\usepackage{tablvar}

\setlength{\parindent}{0mm}
\renewcommand{\arraystretch}{1.5}
\renewcommand{\labelenumi}{\textbf{\theenumi{}.}}
\renewcommand{\labelenumii}{\textbf{\theenumii{}.}}
\setlength{\fboxsep}{3mm}

\setlength{\headheight}{14.5pt}

\setlength{\premulticols}{0pt}

\newcounter{ExoMA}

\newlength{\parindentMA}%
\setlength{\parindentMA}{\parindent}
\addtolength{\parindentMA}{30pt}

\usepackage{simplekv}

\setKVdefault[Theme]{Coopmath,Classiques=false,CANs=false}
\defKV[Theme]{Classique=\setKV[Theme]{Coopmath=false}\setKV[Theme]{Classiques}}
\defKV[Theme]{CAN=\setKV[Theme]{Coopmath=false}\setKV[Theme]{CANs}}

\usepackage{fancyhdr}
\fancyhead[C]{}
\fancyfoot{}

\def\bla{}

\NewDocumentCommand\Theme{o m m m m}{%
  \useKVdefault[Theme]%
  \setKV[Theme]{#1}%
  \ifboolKV[Theme]{CANs}{%
    \usepackage[left=1.5cm,right=1.5cm,top=2cm,bottom=2cm]{geometry}%
    \fancypagestyle{premierePage}{%
      \fancyhead[C]{\textsc{\ifx\bla#2\bla Course aux nombres\else#2\fi}}%
      \fancyfoot[R]{\scriptsize Coopmaths.fr -- CC-BY-SA}%
      \setlength{\headheight}{14.5pt}%
      \NewDocumentCommand\dureeCan{}{\ifx\bla#4\bla 9 minutes\else #4\fi}
      \NewDocumentCommand\titreSujetCan{}{\ifx\bla#3\bla\else\textbf{#3}\fi}
    }%
    \pagestyle{premierePage}
    \colorlet{couleur_theme}{black}%
    \colorlet{couleur_numerotation}{couleur_theme}%
    \let\EXO\EXOCan\let\endEXO\endEXOCan%
  }{%
    \renewcommand{\headrulewidth}{0pt}
    \renewcommand{\footrulewidth}{0pt}
    \pagestyle{fancy}
    \ifboolKV[Theme]{Classiques}{%
      \usepackage[left=0.65cm,right=0.65cm,top=2cm,bottom=1.5cm]{geometry}
      \let\EXO\EXOlibre\let\endEXO\endEXOlibre%
      \colorlet{couleur_theme}{black}%
      \colorlet{couleur_numerotation}{couleur_theme}%
      \fancyhead[C]{\bfseries #3}
      \fancyhead[L]{\bfseries #4}
      \fancyfoot[R]{#5}
    }{%
      \usepackage[left=1.5cm,right=1.5cm,top=4cm,bottom=2cm]{geometry}
      \theme{#2}{#3}{#4}{#5}%
      \let\EXO\EXOcoop\let\endEXO\endEXOcoop%
    }%
  }%
}%

\NewDocumentEnvironment{EXOCan}{m m}{%
}{}

\NewDocumentEnvironment{EXOcoop}{m m}{%
  \begin{tcolorbox}[%
    enhanced,
    mathaleaexo,
    colback=white,
    colframe=white,
    coltitle=black,
    title=\IfNoValueTF{#1}{}{\textmd{#1}},%
    overlay unbroken and first={%
      \IfNoValueTF{#2}{}{%
        \node[
        fill=white,
        anchor=east,
        xshift=10pt,
        text=black
        ]
        at (frame.north east)
        {\scriptsize #2};
      }
      \node[anchor=west,xshift=-20pt,yshift=-15pt] 
      at (frame.north west) {\stepcounter{ExoMA}\numb{\arabic{ExoMA}}};
    }
  ]
}{%
  \end{tcolorbox}
}

\NewDocumentEnvironment{EXOlibre}{m m}{%
  \begin{tcolorbox}[%
    enhanced,
    mathaleaexo,
    colback=white,
    colframe=white,
    coltitle=black,
    title=\stepcounter{ExoMA}\textbf{Exercice \arabic{ExoMA}},%
  ]
  \IfNoValueTF{#1}{}{#1\par}%
}{%
  \end{tcolorbox}%
}


%%% MISE EN PAGE %%%
\usepackage{setspace}
\usepackage{pgf}%
\usepackage{pgfplots}
\pgfplotsset{compat=1.18}
\usetikzlibrary{babel,arrows, arrows.meta,calc,fit,patterns,plotmarks,shapes.geometric,shapes.misc,shapes.symbols,shapes.arrows,shapes.callouts, shapes.multipart, shapes.gates.logic.US,shapes.gates.logic.IEC, er, automata,backgrounds,chains,topaths,trees,petri,mindmap,matrix, calendar, folding,fadings,through,positioning,scopes,decorations.fractals,decorations.shapes,decorations.text,decorations.pathmorphing,decorations.pathreplacing,decorations.footprints,decorations.markings,shadows}

\spaceskip=2\fontdimen2\font plus 3\fontdimen3\font minus3\fontdimen4\font\relax %Pour doubler l'espace entre les mots
\newcommand\numb[1]{ % Dessin autour du numéro d'exercice
  \begin{tikzpicture}[overlay,scale=.8]%,yshift=-.3cm
    \draw[fill=couleur_numerotation,couleur_numerotation](-.4,-0.1)rectangle($(-0.4,.8)+(2em,0)$);
    \draw (-.4,.8) node[white,anchor=north west,inner sep=0.5pt]{\bfseries EX}; 
    \draw[line width=.05cm,couleur_numerotation,fill=white] (0,0)--(.5,.5)--(1,0)--(.5,-.5)--cycle;
    \draw (0.5,0) node[anchor=center,couleur_numerotation]{\large\bfseries#1};
  \end{tikzpicture}
}

% echelle pour le dé
\def\globalscale {0.04}
% abscisse initiale pour les chevrons
\def\xini {3}

\newcommand\theme[4]{%
  \fancyhead[C]{%
    % Tracé du dé
    \begin{tikzpicture}[y=0.80pt, x=0.80pt, yscale=-\globalscale, xscale=\globalscale,remember picture, overlay,transform canvas={xshift=-0.5\linewidth,yshift=2.7cm},fill=couleur_theme]%,xshift=17cm,yshift=9.5cm
      %%%% Arc supérieur gauche%%%%
      \path[fill](523,1424)..controls(474,1413)and(404,1372)..(362,1333)..controls(322,1295)and(313,1272)..(331,1254)..controls(348,1236)and(369,1245)..(410,1283)..controls(458,1328)and(517,1356)..(575,1362)..controls(635,1368)and(646,1375)..(643,1404)..controls(641,1428)and(641,1428)..(596,1430)..controls(571,1431)and(538,1428)..(523,1424)--cycle;
      %%%% Dé face supérieur%%%%
      \path[fill](512,1272)..controls(490,1260)and(195,878)..(195,861)..controls(195,854)and(198,846)..(202,843)..controls(210,838)and(677,772)..(707,772)..controls(720,772)and(737,781)..(753,796)..controls(792,833)and(1057,1179)..(1057,1193)..controls(1057,1200)and(1053,1209)..(1048,1212)..controls(1038,1220)and(590,1283)..(551,1282)..controls(539,1282)and(521,1278)..(512,1272)--cycle;
      %%%% Dé faces gauche et droite%%%%
      \path[fill](1061,1167)..controls(1050,1158)and(978,1068)..(900,967)..controls(792,829)and(756,777)..(753,756)--(748,729)--(724,745)..controls(704,759)and(660,767)..(456,794)..controls(322,813)and(207,825)..(200,822)..controls(193,820)and(187,812)..(187,804)..controls(188,797)and(229,688)..(279,563)..controls(349,390)and(376,331)..(391,320)..controls(406,309)and(462,299)..(649,273)..controls(780,254)and(897,240)..(907,241)..controls(918,243)and(927,249)..(928,256)..controls(930,264)and(912,315)..(889,372)..controls(866,429)and(848,476)..(849,477)..controls(851,479)and(872,432)..(897,373)..controls(936,276)and(942,266)..(960,266)..controls(975,266)and(999,292)..(1089,408)..controls(1281,654)and(1290,666)..(1290,691)..controls(1290,720)and(1104,1175)..(1090,1180)..controls(1085,1182)and (1071,1176)..(1061,1167)--cycle;
      %%%% Arc inférieur bas%%%%
      \path[fill](1329,861)..controls(1316,848)and(1317,844)..(1339,788)..controls(1364,726)and(1367,654)..(1347,591)..controls(1330,539)and(1338,522)..(1375,526)..controls(1395,528)and(1400,533)..(1412,566)..controls(1432,624)and(1426,760)..(1401,821)..controls(1386,861)and(1380,866)..(1361,868)..controls(1348,870)and(1334,866)..(1329,861)--cycle;
      %%%% Arc inférieur gauche%%%%
      \path[fill](196,373)..controls(181,358)and(186,335)..(213,294)..controls(252,237)and(304,190)..(363,161)..controls(435,124)and(472,127)..(472,170)..controls(472,183)and(462,192)..(414,213)..controls(350,243)and(303,283)..(264,343)..controls(239,383)and(216,393)..(196,373)--cycle;
    \end{tikzpicture}
    \begin{tikzpicture}[line cap=round,line join=round,remember picture, overlay, shift={(current page.north west)},yshift=-8.5cm]
      \fill[fill=couleur_theme] (0,5) rectangle (21,6);
      \fill[fill=couleur_theme] (\xini,6)--(\xini+1.5,6)--(\xini+2.5,7)--(\xini+1.5,8)--(\xini,8)--(\xini+1,7)-- cycle;
      \fill[fill=couleur_theme] (\xini+2,6)--(\xini+2.5,6)--(\xini+3.5,7)--(\xini+2.5,8)--(\xini+2,8)--(\xini+3,7)-- cycle;  
      \fill[fill=couleur_theme] (\xini+3,6)--(\xini+3.5,6)--(\xini+4.5,7)--(\xini+3.5,8)--(\xini+3,8)--(\xini+4,7)-- cycle;   
      \node[color=white] at (10.5,5.5) {\LARGE \bfseries{ \MakeUppercase{#4}}};
    \end{tikzpicture}
    \begin{tikzpicture}[remember picture,overlay]
      \node[anchor=north east,inner sep=0pt] at ($(current page.north east)+(0,-.8cm)$) {};
      \node[anchor=west, fill=white, inner sep = 0pt] at ($(current page.north west)+(0.2,-2.3cm)$) {\footnotesize \bfseries{MathALÉA}};
    \end{tikzpicture}
    \begin{tikzpicture}[remember picture,overlay]
      \node[anchor=north east,inner sep=0pt] at ($(current page.north east)+(0,-.8cm)$) {};
      \node[anchor=east, fill=white] at ($(current page.north east)+(-2,-1.5cm)$) {\Huge \textcolor{couleur_theme}{\bfseries{\#}} \bfseries{#2} \textcolor{couleur_theme}{\bfseries \MakeUppercase{#3}}};
    \end{tikzpicture}
  }%
  \fancyfoot[R]{%
    \begin{tikzpicture}[remember picture,overlay]
      \node[anchor=south east] at ($(current page.south east)+(-2,0.25cm)$) {\scriptsize {\bfseries \href{https://coopmaths.fr/}{Coopmaths.fr} -- \href{http://creativecommons.fr/licences/}{CC-BY-SA}}};
    \end{tikzpicture}
    \begin{tikzpicture}[line cap=round,line join=round,remember picture, overlay,xscale=0.5,yscale=0.5, shift={(current page.south west)},xshift=35.7cm,yshift=-6cm]
      \fill[fill=couleur_theme] (\xini,6)--(\xini+1.5,6)--(\xini+2.5,7)--(\xini+1.5,8)--(\xini,8)--(\xini+1,7)-- cycle;
      \fill[fill=couleur_theme] (\xini+2,6)--(\xini+2.5,6)--(\xini+3.5,7)--(\xini+2.5,8)--(\xini+2,8)--(\xini+3,7)-- cycle;  
      \fill[fill=couleur_theme] (\xini+3,6)--(\xini+3.5,6)--(\xini+4.5,7)--(\xini+3.5,8)--(\xini+3,8)--(\xini+4,7)-- cycle;  
    \end{tikzpicture}
  }
  \fancyfoot[C]{}
  \colorlet{couleur_theme}{#1}
  \colorlet{couleur_numerotation}{couleur_theme}
  \def\iconeobjectif{icone-objectif-#1}
  \def\urliconeomethode{icone-methode-#1}
}%

\newcommand*\tikzfootMA{%
  \ifboolKV[Theme]{CANs}{
    \begin{tikzpicture}[remember picture,overlay,shift=(current page.north west)]
      \begin{scope}[x=\textwidth-\oddsidemargin,y=\textheight+5pt]
        \draw[correction,line width=4pt,dashed,dash pattern= on 10pt off 10pt,shift={(-\oddsidemargin,\topmargin)}] (0,0) rectangle (1,-1);
      \end{scope}
    \end{tikzpicture} 
  }{%
  \ifboolKV[Theme]{Classiques}{%
  \begin{tikzpicture}[remember picture,overlay,shift=(current page.south west)]
    \begin{scope}[x={(current page.south east)},y={(current page.north west)}]
      \draw[correction,line width=4pt,dashed,dash pattern= on 10pt off 10pt] ($(0,0)+(5mm,1cm)$) rectangle ($(1,1)+(-5mm,-1.75cm)$);
    \end{scope}
  \end{tikzpicture} 
  }{%
  \begin{tikzpicture}[remember picture,overlay,shift=(current page.south west)]
    \begin{scope}[x={(current page.south east)},y={(current page.north west)}]
      \draw[correction,line width=4pt,dashed,dash pattern= on 10pt off 10pt] ($(0,0)+(5mm,1.5cm)$) rectangle ($(1,1)+(-5mm,-3.75cm)$);
    \end{scope}
  \end{tikzpicture} 
  }%
  }
}

\newenvironment{Correction}{%
  \setcounter{ExoMA}{0}%
  \cfoot{\tikzfootMA}
}{}%

\newcommand\version[1]{
  \fancyhead[R]{
    \begin{tikzpicture}[remember picture,overlay]
    \node[anchor=north east,inner sep=0pt] at ($(current page.north east)+(-.5,-.5cm)$) {\large \textcolor{couleur_theme}{\bfseries V#1}};
    \end{tikzpicture}
  }
}

\usepackage[french]{babel}
\usepackage{icomma}
% loadPackagesFromContent
\usepackage{cancel}
\usepackage{mathrsfs}
\usepackage[tikz]{bclogo}
\newcommand{\e}{\mathrm{e}}
\usepackage{amssymb}
\usepackage{etoolbox}
\newbool{dys}
\setbool{dys}{false}
% ===== POLICES =====
% Le choix est le même en police standard et en police adaptée aux dys ;
% seuls la taille et l'interligne changent entre les deux.
\newcommand{\choiceFonts}[1]{
\ifstrequal{#1}{helvet}{%
  % Helvetica (police standard historique des fiches)
  \usepackage[scaled=1]{helvet}
}{}
\ifstrequal{#1}{Fira}{%
  % Fira Sans + Fira Math
  % Description : Fira Sans est une police moderne, et Fira Math est une version mathématique compatible qui maintient le style sans-serif.
  % Utilisation : Fonctionne bien avec XeLaTeX et LuaLaTeX.
  \usepackage{fontspec}
  \setmainfont{Fira Sans}
  \setsansfont{Fira Sans}
  \usepackage{unicode-math}
  \setmathfont{Fira Math}
}{}
\ifstrequal{#1}{lmodern}{%
  % Latin Modern Sans
  % Description : Une version sans-serif de la célèbre police Latin Modern, qui est compatible avec mathastext.
  % Utilisation : Fonctionne avec pdfLaTeX, XeLaTeX et LuaLaTeX.
  \usepackage{lmodern}
  \renewcommand{\familydefault}{\sfdefault}
  \usepackage{mathastext}
}{}
\ifstrequal{#1}{tgheros}{%
  % TeX Gyre Heros + mathastext
  % Description : TeX Gyre Heros est une alternative à Helvetica, et le package mathastext permet d'utiliser la même police pour les mathématiques.
  % Utilisation : Fonctionne avec pdfLaTeX, XeLaTeX, ou LuaLaTeX.
  \usepackage{tgheros}
  \renewcommand{\familydefault}{\sfdefault}
  \usepackage{mathastext}
}{}
}
\ifbool{dys}{
% POLICE DYS
\choiceFonts{Fira}

%\usepackage{unicode-math}
%\usepackage{fontspec}
%\setmainfont{TeX Gyre Schola}
%\setsansfont{TeX Gyre Schola}
%\setmathfont{TeX Gyre Schola Math}
%\setmainfont{OpenDyslexic}[Scale=1.0]
%\setsansfont{OpenDyslexic}[Scale=1.0]
%\setmathfont{OpenDyslexic}[Scale=1.0,range=up/{Latin,latin,num}]
\usepackage[fontsize=14]{scrextend}
\usepackage{setspace}
\setstretch{1.7}
% Une valeur d'environ 1.2 à 1.7 est souvent recommandée. Cela permet d'augmenter l'espace entre les lignes tout en offrant un léger espacement entre les mots.
\setlength{\spaceskip}{1.2em}
% Une valeur d'environ 1.2em à 1.5em est couramment conseillée. Cela crée un espace plus ample entre les mots, ce qui peut aider à réduire la fatigue visuelle et à améliorer la fluidité de la lecture.
}{
% POLICE STANDARD
%%% EE (24/04/2026) : Cette modif ci-dessous est nécessaire pour accepter ’ comme apostrophe.
% \usepackage[T1]{fontenc}     % Réservé à pdfLaTeX, à remplacer par fontspec en LuaLaTeX
\choiceFonts{helvet}
\usepackage[fontsize=12]{scrextend}
}

\Theme[Classique]{nombres}{}{Fiche calcul 2}{Équations 2}
\begin{document}

\tcbset{mathaleaexo/.style={breakable}}
% @see : https://coopmaths.fr/alea?uuid=71e5c&id=2L21-1&n=2&d=10&s=1&s2=1-2-3-4-5-6-7-8-9&alea=SmM1&cd=1&cols=1
\begin{EXO}{}{}
Justifier si les nombres proposés sont des solutions de l'équation donnée ou non.


\begin{enumerate}
	\item $x^2-14x+45=0~$ pour $~x=5~$, pour $~x=9~$ puis pour $~x=8$
	\item $4x+9=5x-1~$ pour $~x=10~$ puis pour $~x=2$
\end{enumerate}
\tcbset{mathaleaexo/.style={breakable}}
% @see : https://coopmaths.fr/alea?uuid=3ad01&id=2L21-5&n=4&d=10&s=0&s2=false&alea=ML0o&cd=1&cols=1
\end{EXO}

\begin{EXO}{}{}
Résoudre les équations suivantes.

\begin{multicols}{4}
\begin{enumerate}[itemsep=1em]
	\item \begin{minipage}[t]{\linewidth} $8-(2x+8)=2x-4$

 \end{minipage}
	\item \begin{minipage}[t]{\linewidth} $-x-1=-2x+8$

 \end{minipage}
	\item \begin{minipage}[t]{\linewidth} $5(-2x+4)=3x+7$

 \end{minipage}
	\item \begin{minipage}[t]{\linewidth} $x^2+1=(x-8)^2$

 \end{minipage}
\end{enumerate}
\end{multicols}
\tcbset{mathaleaexo/.style={breakable}}
% @see : https://coopmaths.fr/alea?uuid=550cf&id=2L21-6&n=1&d=10&alea=P6aQ&cd=1&cols=1
\end{EXO}

\begin{EXO}{}{}



  Résoudre l'équation $\dfrac{4x+10}{-6}=\dfrac{1}{4}$.
\tcbset{mathaleaexo/.style={breakable}}
% @see : https://coopmaths.fr/alea?uuid=bb6d5&id=2L22-2&n=3&d=10&s=1-2-3-5&alea=Br0N&cd=1&cols=1
\end{EXO}

\begin{EXO}{}{}
Résoudre dans $\mathbb{R}$ :

\begin{multicols}{3}
\begin{enumerate}[itemsep=1em]
	\item \begin{minipage}[t]{\linewidth} $x^{2}-25=0$ 

 \end{minipage}
	\item \begin{minipage}[t]{\linewidth} $x^{2}-32=0$

 \end{minipage}
	\item \begin{minipage}[t]{\linewidth} $4x^{2}-40=0$

 \end{minipage}
\end{enumerate}
\end{multicols}
\tcbset{mathaleaexo/.style={breakable}}
% @see : https://coopmaths.fr/alea?uuid=31017&id=2L22-4&n=1&d=10&s=3&alea=XZGz&cd=1&cols=1
\end{EXO}

\begin{EXO}{}{}



 Préciser les valeurs interdites éventuelles, puis résoudre l'équation : $\dfrac{6+4x}{5x-4}=-1$.

Ensemble des valeurs interdites : 
Ensemble des solutions : 
\tcbset{mathaleaexo/.style={breakable}}
% @see : https://coopmaths.fr/alea?uuid=27741&id=2L20-2&n=1&d=10&s=1&alea=ILxS&cd=1&cols=1
\end{EXO}

\begin{EXO}{}{}
Exprimer le prix total de l'achat, en fonction des lettres introduites dans l'énoncé.


 Nassim veut acheter $2$ poireaux et $3$ navets.

On note $p$ le prix d'un poireau et $n$ le prix d'un navet.
\tcbset{mathaleaexo/.style={breakable}}
% @see : https://coopmaths.fr/alea?uuid=622b9&id=2L20-1&n=1&d=10&s=0&s2=false&alea=lZvf&cd=0&cols=1
\end{EXO}

\begin{EXO}{}{}



 Ursule et Nicolas choisissent un même nombre.

 Ursule lui ajoute $1$ puis multiplie le résultat par $6$ alors que Nicolas lui ajoute $7$ puis multiplie le résultat par $15$.

Ursule et Nicolas obtiennent le même résultat.

Quel nombre commun ont choisi Ursule et Nicolas ?
\tcbset{mathaleaexo/.style={breakable}}
% @see : https://coopmaths.fr/alea?uuid=cf5b7&id=2L20-4&n=1&d=10&s=1&alea=hcGK&cd=1&cols=2
\end{EXO}

\begin{EXO}{}{}


\begin{multicols}{2}
 \begin{minipage}[t]{\linewidth}  En augmentant le côté d'un carré de $5\text{ cm}$, son aire augmente de $76\text{ cm}^2$.

      Quelle est la longueur du côté de ce carré ? 

      On donnera le résultat sous la forme d'une fraction irréductible ou d'un nombre entier le cas échéant. \end{minipage}
\end{multicols}
\end{EXO}


\clearpage

\begin{Correction}
\begin{EXO}{}{}

\begin{enumerate}
	\item Pour $x=5$ : 

$x^2-14\times  x+45=
        5^2-14\times 5+45=
        25-70+45=0$

On trouve bien $0$ pour le membre de gauche donc l'égalité est vraie.

{\color[HTML]{F15929}$x=5$ est donc solution de l'équation $x^2-14x+45=0~$}\medskip

Pour $x=9$ : 

$x^2-14\times  x+45=
        9^2-14\times 9+45=
        81-126+45=0$

On trouve bien $0$ pour le membre de gauche donc l'égalité est vraie.

{\color[HTML]{F15929}$x=9$ est donc solution de l'équation $x^2-14x+45=0~$}\medskip

Pour $x=8$ : 

$x^2-14\times  x+45=
        8^2-14\times 8+45=
        64-112+45=-3$

$-3\not=0$ donc l'égalité n'est pas vraie.

{\color[HTML]{F15929}$x=8$ n'est donc pas solution de l'équation $x^2-14x+45=0~$}\medskip


	\item Pour $x=10$ : 

$4x+9=4\times 10+9=40+9=49$ 

           $5x-1=5\times 10-1=50-1=49$

On trouve le même résultat pour le membre de gauche et pour le membre de droite donc l'égalité est vraie.

{\color[HTML]{F15929}$x=10$ est donc solution de l'équation $4x+9=5x-1~$}\medskip

Pour $x=2$ : 

$4x+9=4\times 2+9=8+9=17$ 

           $5x-1=5\times 2-1=10-1=9$

$17\not=9$ donc l'égalité n'est pas vraie.

{\color[HTML]{F15929}$x=2$ n'est donc pas solution de l'équation $4x+9=5x-1~$}\medskip


\end{enumerate}
\end{EXO}

\begin{EXO}{}{}

\begin{enumerate}[itemsep=1em]
	\item $8-(2x+8)=2x-4$

On développe le membre de gauche.

$8-2x-8=2x-4$

$-2x+0=2x-4$

On soustrait $2x$ aux deux membres.

$-2x+0{\color[HTML]{216D9A}\boldsymbol{-2x}}=2x-4{\color[HTML]{216D9A}\boldsymbol{-2x}}$

$-4x+0=-4$

On ajoute $0$ aux deux membres.

$-4x+0{\color[HTML]{216D9A}\boldsymbol{+0}}=-4{\color[HTML]{216D9A}\boldsymbol{+0}}$

$-4x=-4$

On divise les deux membres par $-4$.

$-4x{\color[HTML]{216D9A}\boldsymbol{\div(-4)}}=-4{\color[HTML]{216D9A}\boldsymbol{\div(-4)}}$

$x=\dfrac{4}{4}$

Par simplification, $x=1$.

La solution de l'équation $8-(2x+8)=2x-4$ est ${\color[HTML]{F15929}\boldsymbol{1}}$.
	\item $-x-1=-2x+8$

On ajoute $2x$ aux deux membres.

$-x-1{\color[HTML]{216D9A}\boldsymbol{+2x}}=-2x+8{\color[HTML]{216D9A}\boldsymbol{+2x}}$

$x-1=8$

On ajoute $1$ aux deux membres.

$x-1{\color[HTML]{216D9A}\boldsymbol{+1}}=8{\color[HTML]{216D9A}\boldsymbol{+1}}$

$x=9$

La solution de l'équation $-x-1=-2x+8$ est ${\color[HTML]{F15929}\boldsymbol{9}}$.
	\item $5(-2x+4)=3x+7$

On développe le membre de gauche.

$-10x+20=3x+7$

On soustrait $3x$ aux deux membres.

$-10x+20{\color[HTML]{216D9A}\boldsymbol{-3x}}=3x+7{\color[HTML]{216D9A}\boldsymbol{-3x}}$

$-13x+20=7$

On soustrait $20$ aux deux membres.

$-13x+20{\color[HTML]{216D9A}\boldsymbol{-20}}=7{\color[HTML]{216D9A}\boldsymbol{-20}}$

$-13x=-13$

On divise les deux membres par $-13$.

$-13x{\color[HTML]{216D9A}\boldsymbol{\div(-13)}}=-13{\color[HTML]{216D9A}\boldsymbol{\div(-13)}}$

$x=\dfrac{13}{13}$

Par simplification, $x=1$.

La solution de l'équation $5(-2x+4)=3x+7$ est ${\color[HTML]{F15929}\boldsymbol{1}}$.
	\item $x^2+1=(x-8)^2$

On développe le membre de droite à l'aide de l'identité remarquable $(a-b)^2=a^2-2ab+b^2$, avec $a=x$ et $b=8$.

$x^2+1=x^2-2\times x\times 8+8^2$

$x^2+1=x^2-16x+64$

On soustrait $x^2$ aux deux membres.

$x^2+1{\color[HTML]{216D9A}\boldsymbol{-x^2}}=x^2-16x+64{\color[HTML]{216D9A}\boldsymbol{-x^2}}$

$1=-16x+64$

On soustrait $64$ aux deux membres.

$1{\color[HTML]{216D9A}\boldsymbol{-64}}=-16x+64{\color[HTML]{216D9A}\boldsymbol{-64}}$

$-63=-16x$

On divise les deux membres par $-16$.

$-63{\color[HTML]{216D9A}\boldsymbol{\div(-16)}}=-16x{\color[HTML]{216D9A}\boldsymbol{\div(-16)}}$

$x=\dfrac{63}{16}$

La solution de l'équation $x^2+1=(x-8)^2$ est ${\color[HTML]{F15929}\boldsymbol{\dfrac{63}{16}}}$.
\end{enumerate}
\end{EXO}

\begin{EXO}{}{}

 
    \begin{bclogo}[couleurBarre=nombres,couleurBord=nombres,epBord=2,couleur=gray!10,logo=\bclampe,arrondi=0.1]{\bf Rappel}
      Pour tout réels $a$, $b$, $c$, $d$ tels que $b$ et $d$ soient non nuls : $\dfrac{a}{b}=\dfrac{c}{d}$ si et seulement si $ad=bc$.
    \end{bclogo}
    \[
\begin{aligned}
\dfrac{4x+10}{-6}&=\dfrac{1}{4}\\
4( 4x+10 )&=-6 \times 1\\
16x+40&=-6\\
16x&=-6-40\\
16x&=-46\\x&=-\dfrac{46}{16}=-\dfrac{23}{8}
\end{aligned}
\]L'équation a donc pour unique solution : ${\color[HTML]{F15929}\boldsymbol{-\dfrac{23}{8}}}$.
\end{EXO}

\begin{EXO}{}{}

\begin{enumerate}[itemsep=1em]
	\item \textbf{Méthode 1 :}

On reconnaît l'identité remarquable $a^2-b^2$ qui se factorise en $(a-b)(a+b)$.

          On l'applique avec $a=x$ et $b=5$.

         Résoudre l'équation revient à résoudre $(x-5)(x+5)=0$ (on reconnaît une équation produit nul).

  
        $x-5=0$ ou $x+5=0$

        $x=5$ ou $x=-5$

Ainsi, $S={\color[HTML]{F15929}\boldsymbol{\{-5;5\}}}$.
            

\textbf{Méthode 2 :}

On isole le carré pour se ramener à une équation du type $x^2=k$.

        Résoudre l'équation revient à résoudre  $x^2=25$.

  
        Puisque $25$ est strictement positif, l'équation a deux solutions : $-\sqrt{25}=-5$ et $\sqrt{25}=5$.

Ainsi, $S={\color[HTML]{F15929}\boldsymbol{\{-5\,;\,5\}}}$.
	\item \textbf{Méthode 1 :}

On reconnaît l'identité remarquable $a^2-b^2$ qui se factorise en $(a-b)(a+b)$.

          On l'applique avec  $ a=x$ et $b=\sqrt{32}$.

         Résoudre l'équation revient à résoudre $(x-\sqrt{32})(x+\sqrt{32})=0$ (on reconnaît une équation produit nul).

  
        $x-\sqrt{32}=0$ ou $x+\sqrt{32}=0$

        $x=\sqrt{32}$ ou $x=-\sqrt{32}$

        Ainsi, $S={\color[HTML]{F15929}\boldsymbol{\{-\sqrt{32}\,;\,\sqrt{32}\}}}$\medskip

\textbf{Méthode 2 :}

On isole le carré pour se ramener à une équation du type $x^2=k$.

        Résoudre l'équation revient à résoudre  $x^2=32$.

  
        Puisque $32$ est strictement positif, l'équation a deux solutions :  $-\sqrt{32}$ et $\sqrt{32}$.

Ainsi, $S={\color[HTML]{F15929}\boldsymbol{\{-\sqrt{32}\,;\,\sqrt{32}\}}}$.
	\item On isole le carré pour se ramener à une équation du type $x^2=k$.

        Résoudre l'équation revient à résoudre  $x^2=10$.

  Puisque $10$  est strictement positif, l'équation a deux solutions :  $-\sqrt{10}$ et $\sqrt{10}$.

Ainsi, $S={\color[HTML]{F15929}\boldsymbol{\{-\sqrt{10}\,;\,\sqrt{10}\}}}$.
\end{enumerate}
\end{EXO}

\begin{EXO}{}{}

 Déterminer les valeurs interdites revient à déterminer les valeurs qui annulent le dénominateur du quotient, puisque la division par $0$ n'existe pas.

Or $5x-4=0$ si et seulement si  $x=\dfrac{4}{5}$. 

          Donc l'ensemble des valeurs interdites est  ${\color[HTML]{F15929}\boldsymbol{\left\{\dfrac{4}{5}\right\}}}$. 

          Pour tout $x\in \mathbb{R}\smallsetminus\left\{\dfrac{4}{5}\right\}$,

            $\begin{aligned}
           \dfrac{6+4x}{5x-4}&=-1\\
            6+4x&=-1\times(5x-4)\,\,\,\,\,\,\,\text{ car les produits en croix sont égaux.}\\
            4x+6&= -5x+4\\
            9x&= -2\\
           x&=-\dfrac{2}{9}
           \end{aligned}$

$-\dfrac{2}{9}$ n'est pas une valeur interdite, donc l'ensemble des solutions de cette équation est  ${\color[HTML]{F15929}\boldsymbol{\mathscr{S}=\left\{-\dfrac{2}{9}\right\}}}$.
\end{EXO}

\begin{EXO}{}{}

 
        Nassim va payer $2$ fois le prix d'un poireau et $3$ fois le prix d'un navet.

 C'est-à-dire $2\times p + 3\times n = 2p + 3n$.

Donc le prix total de l'achat est ${\color[HTML]{F15929}\boldsymbol{2p + 3n}}$.
        
\end{EXO}

\begin{EXO}{}{}

 Posons $x$ le nombre choisi au départ.

Le programme de calcul effectué par Ursule se traduit par : $(x+1)\times 6$.

Le programme de calcul effectué par Nicolas se traduit par : $(x+7)\times 15$.

L'égalité des résultats se traduit par l'équation suivante :

$(x+1)\times 6=(x+7)\times 15$

{\bfseries \color{black}Résolvons cette équation :}

  Commençons par développer et réduire les deux membres :

  $\begin{aligned}(x + 1)\times 6 &= 6x + 6\\
  (x + 7)\times 15 &= 15x + 105\\
  \end{aligned}$

  Ainsi, l'équation devient : $6x + 6 = 15x + 105$.

  
  Soustrayons $15x$ des deux membres et soustrayons $6$ des deux membres :

  $\begin{aligned}6x - 15x &= 105 - 6\\
  -9\times x &= 99\\
  \end{aligned}$.

  Divisons les deux membres par $-9$ :

  $x = \dfrac{99}{-9} = -11$.

{\bfseries \color{black}Vérification :}

 $(-11 + 1)\times 6 = -10\times 6 = -60$

 
  $(-11 + 7)\times 15 = -4\times 15 = -60$\medskip

Ursule et Nicolas ont donc choisi au départ le nombre ${\color[HTML]{F15929}\boldsymbol{-11}}$.
\end{EXO}

\begin{EXO}{}{}
\begin{multicols}{2}
 \begin{minipage}[t]{\linewidth} On note $x$ la longueur du côté du carré que l'on cherche.

      La différence des aires entre les deux carrés est  $76$, on cherche $x$ tel que :   

      $\begin{aligned}
      (x+5)^2-x^2&=76\\
      \cancel{x^2}+ 10x+5^2-\cancel{x^2}&=76\\
       10x+25&=76\\
        10x+25{\color[HTML]{F15929}\boldsymbol{-25}}&=76{\color[HTML]{F15929}\boldsymbol{-25}}\\
         10x&=51\\
         \dfrac{10x}{{\color[HTML]{F15929}\boldsymbol{10}}}&=\dfrac{51}{{\color[HTML]{F15929}\boldsymbol{10}}}\\
         x&=\dfrac{51}{10}\end{aligned}$
           

           La longueur du côté du carré est  $\dfrac{51}{10}\text{ cm}$.
  \end{minipage}
\end{multicols}
\end{EXO}

\clearpage
\end{Correction}
\end{document}

% Local Variables:
% TeX-engine: luatex
% End: